%% GENERATED by tools/gen_latex.py -- edit content/ch11.py instead.
%% Build:  latexmk -pdf ch11.tex     (or: pdflatex ch11.tex, twice)
\documentclass[aspectratio=169,11pt,t]{beamer}
\usepackage{itbpro}

\coursetitle{Image Segmentation}
\coursesubtitle{A class for every pixel --- built from scratch as an encoder-decoder, then handed to a pretrained model you prompt rather than train.}
\coursekicker{Chapter 11}
\coursesource{Chollet \& Watson, Deep Learning with Python 3e --- chapter 11}
\courseduration{2.5 hours}
\coursepresenter{Rahman Indra Kesuma, S.Kom., M.Cs. $\cdot$ Prof. Bambang Riyanto Trilaksono}{Teaching Assistant $\cdot$ Lead Instructor}
\coursebrand{AI for Professional \textperiodcentered\ ITB \textperiodcentered\ Ch 11}

\title{Image Segmentation}
\author{Rahman Indra Kesuma, S.Kom., M.Cs. $\cdot$ Prof. Bambang Riyanto Trilaksono}
\date{}

\begin{document}
\covertitle

\framekicker{Learning outcomes}
\begin{frame}[shrink=25]{Session Objectives}
\leadpar{By the end of this session, participants can:}
\begin{isteps}
\item Place \textbf{classification, segmentation, and detection} --- the three tasks almost all computer vision reduces to.
\item Distinguish \textbf{semantic, instance, and panoptic} segmentation.
\item Build an \textbf{encoder-decoder} segmentation model, and explain why the output must match the input's spatial size.
\item Say why segmentation uses \textbf{strided convolutions instead of max pooling}.
\item Explain what \textbf{\inlinecode{Conv2DTranspose}} does and how it undoes a strided convolution.
\item Compute and interpret \textbf{Intersection over Union}, and configure the Keras metric correctly for sparse targets.
\item Use \textbf{Segment Anything} with point and box prompts, without fine-tuning anything.
\end{isteps}

\vskip 6pt
\vskip 4pt
\reslink{Book}{Chapter 11 --- full text}{https://deeplearningwithpython.io/chapters/chapter11\_image-segmentation}
\reslink{Notebook}{01\_segmentation\_from\_scratch.ipynb}{../../notebooks/ch11/01\_segmentation\_from\_scratch.ipynb}
\reslink{Notebook}{02\_conv2dtranspose.ipynb}{../../notebooks/ch11/02\_conv2dtranspose.ipynb}
\reslink{Notebook}{03\_segment\_anything.ipynb}{../../notebooks/ch11/03\_segment\_anything.ipynb}
\reslink{Repo}{Author's official notebook}{https://github.com/fchollet/deep-learning-with-python-notebooks/blob/master/chapter11\_image-segmentation.ipynb}
\vskip 2pt
\end{frame}

\framekicker{Section 11.1}
\begin{frame}[shrink=25]{Three tasks almost everything reduces to}
\begin{center}
\includegraphics[width=0.92\linewidth,height=0.60\textheight,keepaspectratio]{figs/ch11-tasks.tex.pdf}
\end{center}
\vskip 3pt{\color{ink3}\fontsize{7.4}{9.4}\selectfont Figure 11.1 --- classification, segmentation, and detection.\par}
\end{frame}

\framekicker{Section 11.1}
\begin{frame}[shrink=25]{\ldots{}and the more specialised ones beyond them}
{\fontsize{9.4}{12.6}\selectfont\color{fg2}Computer vision also covers \textbf{image similarity scoring, keypoint detection} (facial features, for instance), \textbf{pose estimation, 3D mesh estimation}, and \textbf{depth estimation}.\par}\vskip 4pt
\begin{iband}[signal]
But classification, segmentation, and detection are the foundation every engineer should know: \hilite{almost all computer vision applications boil down to one of these three}
\end{iband}
\end{frame}

\framekicker{Section 11.1.1}
\begin{frame}[shrink=25]{Three flavours of segmentation}
\begin{center}
\includegraphics[width=0.92\linewidth,height=0.60\textheight,keepaspectratio]{figs/ch11-flavours.tex.pdf}
\end{center}
\vskip 3pt{\color{ink3}\fontsize{7.4}{9.4}\selectfont Figure 11.2 --- panoptic is the most informative of the three.\par}
{\fontsize{9.4}{12.6}\selectfont\color{fg2}All of them assign a class \textbf{to each pixel}, partitioning the image into zones --- \textit{background} and \textit{foreground}, or \textit{road}, \textit{car}, and \textit{sidewalk}.\par}\vskip 4pt
\end{frame}

\framekicker{Section 11.1.1}
\begin{frame}[shrink=25]{Where it is actually used}
\begin{minipage}[t]{0.241\textwidth}
\begin{icard}
\cardh{Image and video editing}
Background replacement, object removal, compositing.
\end{icard}
\end{minipage}%
\hfill
\begin{minipage}[t]{0.241\textwidth}
\begin{icard}
\cardh{Autonomous driving}
Road, lane, and obstacle segmentation at pixel precision.
\end{icard}
\end{minipage}%
\hfill
\begin{minipage}[t]{0.241\textwidth}
\begin{icard}
\cardh{Robotics}
Knowing exactly where a graspable object begins and ends.
\end{icard}
\end{minipage}%
\hfill
\begin{minipage}[t]{0.241\textwidth}
\begin{icard}
\cardh{Medical imaging}
Delineating an organ or a lesion rather than merely detecting it.
\end{icard}
\end{minipage}%
\end{frame}

\coursesection{01}{Training a segmentation model from scratch}{An encoder that compresses, and a decoder that puts it back.}

\framekicker{Section 11.2.1}
\begin{frame}[shrink=25]{What a label looks like when the label is an image}
\begin{center}
\includegraphics[width=0.92\linewidth,height=0.60\textheight,keepaspectratio]{figs/ch11-mask.tex.pdf}
\end{center}
\vskip 3pt{\color{ink3}\fontsize{7.4}{9.4}\selectfont A segmentation mask is the equivalent of a label: same size as the input, one channel, one integer per pixel.\par}
{\fontsize{9.4}{12.6}\selectfont\color{fg2}The dataset is \textbf{Oxford-IIIT Pets}: 7,390 pictures of cat and dog breeds, each with a foreground-background mask.\par}\vskip 4pt
\end{frame}

\framekicker{Section 11.2.1}
\begin{frame}[shrink=25]{Pairing images with their masks}
{\fontsize{9.4}{12.6}\selectfont\color{fg2}Inputs are JPGs in \inlinecode{images/}; each mask is a PNG of the same name in \inlinecode{annotations/trimaps/}. Sorting both lists is what keeps them aligned.\par}\vskip 4pt
\icode{python}{building the two path lists}{listings/ch11-01.py}
\ioutput{listings/ch11-02.txt}
\begin{iband}[amber]
The \inlinecode{[!.]} in the glob is not cosmetic: the trimaps directory contains hidden files, and including them would \hilite{silently offset every image-mask pair after the first one}
\end{iband}
\end{frame}

\framekicker{Section 11.2.1}
\begin{frame}[shrink=25]{Shuffling two lists in lockstep}
{\fontsize{9.4}{12.6}\selectfont\color{fg2}The paths arrive sorted by breed, so they must be shuffled --- but images and masks have to stay paired. The trick is the same seed, twice.\par}\vskip 4pt
\icode{python}{loading into memory}{listings/ch11-03.py}
\begin{iband}[signal]
Two details worth naming: the labels are shifted from \textbf{1, 2, 3 to 0, 1, 2} because the loss expects zero-based classes, and the whole dataset is loaded into memory because \hilite{7,390 images at 200$\times$200 comfortably fits}
\end{iband}
\end{frame}

\framekicker{Section 11.2.1}
\begin{frame}[shrink=25]{Looking at a mask before training on it}
{\fontsize{9.4}{12.6}\selectfont\color{fg2}Chapter 6 insisted on looking at your data. For masks that needs a small transformation, because values of 1, 2 and 3 are all indistinguishable black on screen.\par}\vskip 4pt
\icode{python}{making a mask visible}{listings/ch11-04.py}
{\fontsize{9.4}{12.6}\selectfont\color{fg2}The result shows the animal in black, the background in near-white, and a grey contour band between them --- \hilite{the third class, which is easy to forget exists}\par}\vskip 4pt
\end{frame}

\framekicker{Section 11.2.2}
\begin{frame}[shrink=25]{The shape of the problem}
\begin{center}
\includegraphics[width=0.92\linewidth,height=0.60\textheight,keepaspectratio]{figs/ch11-encdec.tex.pdf}
\end{center}
\vskip 3pt{\color{ink3}\fontsize{7.4}{9.4}\selectfont Down three times by a factor of two, then back up three times.\par}
{\fontsize{9.4}{12.6}\selectfont\color{fg2}The first half is an ordinary classification-style ConvNet: it \textbf{compresses} the image into a small feature map where each location carries information about a large chunk of the original.\par}\vskip 4pt
\end{frame}

\framekicker{Section 11.2.2}
\begin{frame}[shrink=25]{The model, in one function}
{\fontsize{9.4}{12.6}\selectfont\color{fg2}Note two things: \inlinecode{padding="same"} everywhere, and the absence of any pooling layer.\par}\vskip 4pt
\icode{python}{the encoder half}{listings/ch11-05.py}
{\fontsize{9.4}{12.6}\selectfont\color{fg2}Familiar so far --- growing filter counts, shrinking spatial size. \textbf{Except} that every downsample is a stride rather than a pooling layer.\par}\vskip 4pt
\end{frame}

\framekicker{Section 11.2.2}
\begin{frame}[shrink=25]{\ldots{}and the decoder half}
{\fontsize{9.4}{12.6}\selectfont\color{fg2}The second half mirrors the first, reversing each stride to climb back to the input's spatial size.\par}\vskip 4pt
\icode{python}{the decoder half}{listings/ch11-06.py}
{\fontsize{9.4}{12.6}\selectfont\color{fg2}Three strided convolutions down, three strided transposed convolutions up --- \hilite{the counts have to match} or the output will not align with the mask.\par}\vskip 4pt
\end{frame}

\framekicker{Section 11.2.2}
\begin{frame}[shrink=25]{What that last layer is doing}
\begin{iband}[signal]
The final layer is a \textbf{per-pixel softmax}: for every one of the 200$\times$200 output positions it produces a distribution over the three classes. \hilite{Classification, repeated 40,000 times.}
\end{iband}
{\fontsize{9.4}{12.6}\selectfont\color{fg2}Which is why \inlinecode{padding="same"} appears on every layer: any border trimming would leave the output a different size from the mask it is being compared against.\par}\vskip 4pt
\end{frame}

\framekicker{Section 11.2.2}
\begin{frame}[shrink=25]{Why no max pooling here}
\begin{center}
\includegraphics[width=0.92\linewidth,height=0.60\textheight,keepaspectratio]{figs/ch11-poolvstride.tex.pdf}
\end{center}
\vskip 3pt{\color{ink3}\fontsize{7.4}{9.4}\selectfont The trade-off that was invisible in chapter 8 becomes decisive here.\par}
\end{frame}

\framekicker{Section 11.2.2}
\begin{frame}[shrink=25]{The rule that follows from it}
{\fontsize{9.4}{12.6}\selectfont\color{fg2}2$\times$2 max pooling returns one scalar per window with \textbf{zero knowledge of which of the four positions it came from}. For classification that loss is harmless. For segmentation it is fatal --- the output \textit{is} a map of locations.\par}\vskip 4pt
\begin{iband}[signal]
So the general rule the book states: \textbf{use strides instead of max pooling in any model that cares about feature location} --- including \hilite{the generative models of chapter 17}
\end{iband}
\end{frame}

\framekicker{Section 11.2.2}
\begin{frame}[shrink=25]{Conv2DTranspose: a convolution that learns to upsample}
{\fontsize{9.4}{12.6}\selectfont\color{fg2}The encoder leaves a \inlinecode{(25, 25, 256)} feature map, but the output must be \inlinecode{(200, 200, 3)}. Something has to reverse the downsampling.\par}\vskip 4pt
\icode{python}{the two are inverses}{listings/ch11-07.py}
\begin{iband}[signal]
So a stack of \inlinecode{Conv2D} with strides, mirrored by a stack of \inlinecode{Conv2DTranspose} with the same strides, \hilite{returns you to the original spatial size} --- which is exactly what the model above does.
\end{iband}
\end{frame}

\framekicker{Section 11.2.2}
\begin{frame}[shrink=25]{Measuring it: Intersection over Union}
\begin{center}
\includegraphics[width=0.92\linewidth,height=0.60\textheight,keepaspectratio]{figs/ch11-iou.tex.pdf}
\end{center}
\vskip 3pt{\color{ink3}\fontsize{7.4}{9.4}\selectfont Accuracy per pixel would be dominated by the background; IoU is not.\par}
{\fontsize{9.4}{12.6}\selectfont\color{fg2}It can be computed per class or averaged across classes, and it is the standard measure of how well a predicted mask matches the ground truth.\par}\vskip 4pt
\end{frame}

\framekicker{Section 11.2.2}
\begin{frame}[shrink=25]{Configuring the Keras metric correctly}
{\fontsize{9.4}{12.6}\selectfont\color{fg2}Keras has IoU built in, but two arguments have to match your data --- and getting them wrong produces a silently meaningless number.\par}\vskip 4pt
\icode{python}{the IoU metric}{listings/ch11-08.py}
\begin{iband}[amber]
\inlinecode{sparse\_y\_true=True} with \inlinecode{sparse\_y\_pred=False} is the combination you almost always want here: \textbf{targets are integers, predictions are probabilities}. \hilite{Mismatching them does not raise an error.}
\end{iband}
\end{frame}

\framekicker{Section 11.2.2}
\begin{frame}[shrink=25]{Compiling and fitting it}
{\fontsize{9.4}{12.6}\selectfont\color{fg2}From here it is the ordinary workflow: sparse categorical crossentropy, because the targets are integer class IDs per pixel.\par}\vskip 4pt
\icode{python}{training the segmentation model}{listings/ch11-09.py}
{\fontsize{9.4}{12.6}\selectfont\color{fg2}Everything you learned in chapters 6 and 7 applies unchanged --- only the model's output shape and the metric are new.\par}\vskip 4pt
\end{frame}

\framekicker{Section 11.2.2}
\begin{frame}[shrink=25]{Splitting, and what the arrays look like}
{\fontsize{9.4}{12.6}\selectfont\color{fg2}A thousand samples are reserved for validation. Because both arrays were shuffled with the same seed, a plain slice is enough.\par}\vskip 4pt
\icode{python}{the split}{listings/ch11-10.py}
\ioutput{listings/ch11-11.txt}
\begin{iband}[signal]
Note the target shape: \textbf{one channel, not three}. The input is RGB; the target is \hilite{a single integer per pixel}
\end{iband}
\end{frame}

\framekicker{Section 11.2.2}
\begin{frame}[shrink=25]{Predicting, and reading the output}
{\fontsize{9.4}{12.6}\selectfont\color{fg2}The model outputs probabilities per class per pixel, so turning that into a picture means taking an argmax along the channel axis.\par}\vskip 4pt
\icode{python}{predicting a mask}{listings/ch11-12.py}
\ioutput{listings/ch11-13.txt}
{\fontsize{9.4}{12.6}\selectfont\color{fg2}That \inlinecode{argmax} is the segmentation equivalent of the \inlinecode{predictions[0].argmax()} from chapter 2 --- \hilite{performed once per pixel instead of once per image}\par}\vskip 4pt
\end{frame}

\coursesection{02}{Using a pretrained segmentation model}{Segment Anything: prompt it, do not train it.}

\framekicker{Section 11.3}
\begin{frame}[shrink=25]{What SAM is}
{\fontsize{9.4}{12.6}\selectfont\color{fg2}The \textbf{Segment Anything Model}, from Meta AI, released April 2023. Trained on \textbf{11 million images} and their masks, covering \textbf{over 1 billion object instances}.\par}\vskip 4pt
\begin{iband}[signal]
The main innovation: it is \textbf{not limited to a predefined set of object classes}. You segment a new kind of object simply by giving an example of what you are looking for --- \hilite{and you do not need to fine-tune it first}
\end{iband}
\end{frame}

\framekicker{Section 11.3.2}
\begin{frame}[shrink=25]{The dataset and the model were built together}
\begin{center}
\includegraphics[width=0.92\linewidth,height=0.60\textheight,keepaspectratio]{figs/ch11-samloop.tex.pdf}
\end{center}
\vskip 3pt{\color{ink3}\fontsize{7.4}{9.4}\selectfont The partially trained model was used to help label the data it would later be trained on.\par}
{\fontsize{9.4}{12.6}\selectfont\color{fg2}The goal of SA-1B is \textbf{fully segmented images}: every object given a unique mask. Each image carries around \textbf{100 masks on average}, and some have over \textbf{500}.\par}\vskip 4pt
\end{frame}

\framekicker{Section 11.3.2}
\begin{frame}[shrink=25]{Three components, trained together}
\begin{center}
\includegraphics[width=0.92\linewidth,height=0.60\textheight,keepaspectratio]{figs/ch11-sam.tex.pdf}
\end{center}
\vskip 3pt{\color{ink3}\fontsize{7.4}{9.4}\selectfont Figure 11.8 --- image encoder, prompt encoder, mask decoder. The decoder emits several candidate masks, each with a score.\par}
{\fontsize{9.4}{12.6}\selectfont\color{fg2}The image encoder is \textbf{something you already know how to build} --- it is much like the Xception of chapters 8 and 9. The prompt encoder and mask decoder use techniques from later chapters.\par}\vskip 4pt
\end{frame}

\framekicker{Section 11.3.2}
\begin{frame}[shrink=25]{How the training triples are generated}
{\fontsize{9.4}{12.6}\selectfont\color{fg2}The model is trained on \textbf{(image, prompt, mask)} triples, and the way they are produced is strikingly simple.\par}\vskip 4pt
\begin{isteps}
\item For a given image, \textbf{choose a random mask} from its annotations.
\item \textbf{Randomly choose} whether to create a box prompt or a point prompt.
\item For a \textbf{point prompt}, pick a random pixel inside the mask.
\item For a \textbf{box prompt}, draw a box around all points inside the mask.
\end{isteps}
\begin{iband}[signal]
That can be repeated indefinitely, sampling \textbf{many triples from each image} --- which is how 11 million images become \hilite{a far larger number of training examples}
\end{iband}
\end{frame}

\framekicker{Section 11.3.4}
\begin{frame}[shrink=25]{Two kinds of prompt}
\begin{minipage}[t]{0.494\textwidth}
\begin{icardaccent}
\cardh{Point prompts}
Select a point and let the model segment the object it belongs to. Points are \textbf{labelled}: 1 for foreground, 0 for background.
\end{icardaccent}
\end{minipage}%
\hfill
\begin{minipage}[t]{0.494\textwidth}
\begin{icardaccent}
\cardh{Box prompts}
Draw an approximate box around an object --- \textbf{it does not need to be precise} --- and let the model segment what is inside.
\end{icardaccent}
\end{minipage}%
\begin{iband}[signal]
In ambiguous cases, pass \textbf{several labelled points} rather than one: points labelled 1 say what to include, points labelled 0 say \hilite{what to leave out}
\end{iband}
\end{frame}

\framekicker{Section 11.3.4}
\begin{frame}[shrink=25]{Prompting it with a single point}
{\fontsize{9.4}{12.6}\selectfont\color{fg2}The book uses a photograph of a bowl of fruit and asks for whatever object lies under one chosen pixel.\par}\vskip 4pt
\icode{python}{a point prompt}{listings/ch11-14.py}
\begin{iband}[signal]
One point, no training, no class list --- and SAM returns a mask for the peach that point landed on. \hilite{That is the shift this section is about.}
\end{iband}
\end{frame}

\framekicker{Section 11.3.4}
\begin{frame}[shrink=25]{Reading the output}
{\fontsize{9.4}{12.6}\selectfont\color{fg2}The decoder returns \textbf{several candidate masks with scores}, because a point is genuinely ambiguous: it might mean the peach, the bowl, or the whole arrangement.\par}\vskip 4pt
\icode{python}{picking a mask}{listings/ch11-15.py}
\begin{iband}[amber]
Taking the top-scoring mask is the default, but the ambiguity is real information: \textbf{if the three candidates disagree wildly, the prompt itself was ambiguous} --- and \hilite{an extra background point usually resolves it}
\end{iband}
\end{frame}

\framekicker{Section 11.3}
\begin{frame}[shrink=25]{From scratch, or pretrained?}
\begin{itable}{>{\raggedright\arraybackslash}p{0.2068\textwidth}>{\raggedright\arraybackslash}p{0.3666\textwidth}>{\raggedright\arraybackslash}p{0.3666\textwidth}}
\thead{} & \thead{From scratch (§11.2)} & \thead{Segment Anything (§11.3)} \\ \midrule
\textbf{What you need} & A labelled mask for every training image. & One point or box, at inference time. \\
\textbf{Classes} & Fixed at training time. & \textbf{None} --- it segments whatever you point at. \\
\textbf{Training cost} & A full training run. & Zero. \\
\textbf{When it wins} & You need a \textbf{specific, repeatable} class set --- medical structures, industrial defects --- and you have the masks. & You need \textbf{generic} object masks, or you have no labelled data at all. \\
\end{itable}
\end{frame}

\framekicker{Section 11.3.1}
\begin{frame}[shrink=25]{Loading SAM}
{\fontsize{9.4}{12.6}\selectfont\color{fg2}Like the classifiers in chapter 8, it comes from KerasHub with a matching preprocessor --- the model expects images at a particular size and range.\par}\vskip 4pt
\icode{python}{instantiating Segment Anything}{listings/ch11-16.py}
{\fontsize{9.4}{12.6}\selectfont\color{fg2}No training data, no class list, no fine-tuning step --- \hilite{the next thing that happens is a prediction}\par}\vskip 4pt
\end{frame}

\framekicker{Section 11.3}
\begin{frame}[shrink=25]{Why a promptable model changes the workflow}
\begin{minipage}[t]{0.494\textwidth}
\begin{icardwarn}
\cardh{The old shape}
Decide your classes $\rightarrow$ collect masks for each $\rightarrow$ train $\rightarrow$ the model can segment \textbf{those classes and no others}.
\end{icardwarn}
\end{minipage}%
\hfill
\begin{minipage}[t]{0.494\textwidth}
\begin{icardgood}
\cardh{The new shape}
Point at the thing $\rightarrow$ get a mask. The class set is \textbf{decided at inference time}, by the person using it.
\end{icardgood}
\end{minipage}%
\begin{iband}[signal]
That is the same move as prompting a language model, arriving in vision --- and it is why chapter 1 grouped \textbf{foundation models} as a category rather than a technique.
\end{iband}
\end{frame}

\framekicker{Section 11.3.4}
\begin{frame}[shrink=25]{Refining an ambiguous prompt}
{\fontsize{9.4}{12.6}\selectfont\color{fg2}A single point on a peach might mean the peach, the bowl, or the whole still life. Extra labelled points resolve it.\par}\vskip 4pt
\icode{python}{adding a background point}{listings/ch11-17.py}
\begin{iband}[signal]
This is the practical loop with SAM: \textbf{prompt, look, add a point, prompt again}. It is closer to \hilite{using a selection tool} than to training a model.
\end{iband}
\end{frame}

\framekicker{Practice}
\begin{frame}[shrink=25]{Where segmentation quietly goes wrong}
\begin{minipage}[t]{0.494\textwidth}
\begin{icardbad}
\cardh{Images and masks drift apart}
Two lists sorted or shuffled differently. Nothing errors --- the model simply learns nothing, and the loss curve looks like a capacity problem.
\end{icardbad}
\end{minipage}%
\hfill
\begin{minipage}[t]{0.494\textwidth}
\begin{icardbad}
\cardh{Off-by-one labels}
Masks numbered 1..N fed to a loss expecting 0..N-1. The last class is never predicted and the first is never learned.
\end{icardbad}
\end{minipage}%


\vskip 4pt

\begin{minipage}[t]{0.494\textwidth}
\begin{icardwarn}
\cardh{Accuracy instead of IoU}
If 90\% of pixels are background, a model predicting \textit{background} everywhere scores 90\% accuracy and \textbf{zero} foreground IoU.
\end{icardwarn}
\end{minipage}%
\hfill
\begin{minipage}[t]{0.494\textwidth}
\begin{icardwarn}
\cardh{Padding mismatch}
One layer without \inlinecode{padding="same"} and the output no longer matches the mask's size.
\end{icardwarn}
\end{minipage}%
\end{frame}
\note{All four show up in the lab. The first two are the ones that look like modelling problems and are not.}

\framekicker{Summary}
\begin{frame}[shrink=25]{What has to survive this chapter}
\begin{isteps}
\item \textbf{Classification, segmentation, detection} --- almost every vision application is one of these three.
\item \textbf{Semantic, instance, panoptic} --- increasing amounts of information per pixel.
\item A segmentation model is an \textbf{encoder-decoder}: compress, then restore the original spatial size.
\item \textbf{Strides, not max pooling}, whenever location matters --- pooling destroys position within the window.
\item \textbf{\inlinecode{Conv2DTranspose} learns to upsample}, undoing a strided convolution.
\item \textbf{IoU, not accuracy} --- and check \inlinecode{sparse\_y\_true} / \inlinecode{sparse\_y\_pred} against your data.
\end{isteps}
\end{frame}

\framekicker{Summary}
\begin{frame}[shrink=25]{\ldots{}and the shift that SAM represents}
\begin{iband}[signal]
A pretrained model that is \textbf{prompted rather than trained}, with no predefined class list. It is the first model in this book that works that way --- and \hilite{the pattern returns in every chapter from 15 onwards}
\end{iband}
\vskip 4pt
\reslink{NOTEBOOK}{03\_segment\_anything.ipynb}{../../course-slides/notebooks/ch11/03\_segment\_anything.ipynb}
\reslink{NEXT}{Chapter 12 --- Object detection}{../ch12/index.html}
\vskip 2pt
\end{frame}

\end{document}
